\scanpage{48}
\mspage{4)}

On a un foncteur
\[
 \Gamma:\quad S_{*}\unsure{\mathscr T|\mathbb T}{\unsure{\mathbb D|D}}_{\mathrm{top}}^{\pm}
 \longrightarrow\deleted{\unsure{\Gamma MD}}(\mathrm{Grabipond})\tag*{$(*)$}
\]
qui associe à un objet \inserted{$(X,\widetilde\Sigma,R)$} de $S\mathscr T\deleted{\unsure{D}}\mathbb D_{\mathrm{top}}^{\pm}$ le graphe \inserted{$MD$} $(S,\vec A,\sigma,\deleted{\unsure{r|k}}o,\deleted{\unsure{g}}\gamma,\rho)$ suivant :

\[
 S=\pi_0(X\deleted{\unsure{,}}\setminus\widetilde\Sigma)\xleftarrow{\sim}\pi_0(\mathrm{Déc}(X,\widetilde\Sigma))
\]

$\vec A={}$ ens des couples $(C,\deleted{\unsure{g}}\omega)$, où $C$ \deleted{est un} \inserted{$\in\pi_0(\widetilde\Sigma)$}, et où $\omega\in\mathrm{Or}(C)$ \inserted{$\simeq\mathrm{rive}_{X}(C\deleted{\unsure{,}})$}, soumis à la condition que si $C\in\pi_0(\partial X)$, alors $\omega$ soit l'orientation induite i.e. la rive qui se trouve dans $X$

$\sigma={}$ échange des rives (quand on a le droit d'échanger, sinon, i.e. si \unsure{$^{a}$}$C\in\pi_0(\partial X)$, on pose $\sigma(a)=a$)

\deleted{$\unsure{R\to R}$} $o\deleted{\unsure{(a)}}:$ associe à $(C,\text{rive de }C)$ l'unique ``face'' $\in\pi_0(X\setminus\widetilde\Sigma)$ qui contient cette rive

\[
 \gamma:\underset{\substack{\shortparallel\\\pi_0(X\setminus\widetilde\Sigma)}}{S}\longrightarrow\mathbb N
 \qquad\text{application \uline{genre} du composant $\unsure{X'|X_i|X_s}$.}
\]

\[
 \deleted{\unsure{\rho}}r:\quad A=\pi_0(\widetilde\Sigma)\longrightarrow\mathbb N:\quad C\longmapsto\mathrm{Card}(R\cap C).
\]

Le foncteur $(*)$ \deleted{est \uline{ess. surjectif} et} \inserted{induit un \uline{iso} sur les} $\pi_0$ et \unsure{une|un} surjectif sur les $\pi_1$ (c'est un \uline{$0$-isomorphisme} \deleted{\unsure{)}}, en langage homotopique).

\deleted{\unsure{Prop.}} \deleted{\unsure{Les obj.}}

\uline{Remarques} \quad Dans $S_{*}\mathscr T\mathbb D_{\mathrm{top}}^{\pm}$ il y a une opération évidente de \uline{sommes} \inserted{\uline{directes}}, dans \deleted{$\Gamma MD$} $\mathrm{Grabipond}$ aussi, et $(*)$ y commute.
% archive sheet number: 47

